% -----------------------------------------------------------------------------
% 标量/赝标量 Yukawa 散射：主文件保留目标式与物理主线
% -----------------------------------------------------------------------------

局域接触相互作用提供一个没有传播子角结构的参照。对
\begin{equation*}
 \Lag_{\rm int}=-\frac{g_4}{4!\hbar c}\varphi^4,
 \qquad [g_4]=1,
\end{equation*}
Wick 收缩给出树级约化顶角 $-\ii g_4$，因此
\begin{equation}
 \ii\mathcal M=-\ii g_4,
 \qquad
 \overline{|\mathcal M|^2}=g_4^2.
 \label{eq:phi4-M-tree-si-r10}
\end{equation}
四个外粒子质量相同且末态两个标量全同，质心系微分截面为
\begin{equation}
 \boxed{
 \frac{\dd\sigma}{\dd\Omega}
 =\frac{\hbar^2g_4^2}{128\pi^2s}.}
 \label{eq:phi4-isotropic-cross-si-r10}
\end{equation}
局域、无导数的接触顶角不携带方向信息，所以该角分布各向同性。非平凡角结构必须来自传播子分母或自旋、极化分子。

接触 $\varphi^4$ 理论只作为“无传播子角结构”的比较基准；有限质量媒介交换必须从自己的局域场论独立建立。现在引入一个 Dirac 场与实标量场的局域耦合，采用全书统一的 SI 场归一化，
\begin{equation}
 \boxed{
 \Lag_Y
 =\bar\psi(\ii\hbar c\slashed\partial-mc^2)\psi
 +\frac12(\partial_\mu\phi)(\partial^\mu\phi)
 -\frac12\frac{m_Y^2c^2}{\hbar^2}\phi^2
 -y\sqrt{\hbar c}\,\phi\bar\psi\psi.}
 \label{eq:yukawa-full-si-r10}
\end{equation}
其中 $m$ 是费米子质量，$m_Y$ 是标量媒介粒子的静质量，$y$ 是无量纲的标量--费米子三线性耦合常数；这种局域三线性耦合通常称为 Yukawa 耦合。由自由二次核和相互作用项分别得到约化传播子与顶角
\begin{equation}
 S_{\rm red}(p)=\frac{\ii(\slashed p+mc)}{p^2-m^2c^2+\ii0^+},
 \qquad
 \Delta_{\rm red}(q)=\frac{\ii}{q^2-m_Y^2c^2+\ii0^+},
 \qquad
 V_Y=-\ii y.
 \label{eq:yukawa-propagators-si-r10}
\end{equation}
这里四动量采用 $p^\mu=(E/c,\mathbf p)$，因而 $p^2$ 与 $m^2c^2$ 都具有动量平方量纲。

对可区分的两种费米子 $A,B$，种类守恒使最低阶只有 $t$ 通道。令 $t=(p_1-p_3)^2$，树级约化振幅为
\begin{equation}
 \boxed{
 \mathcal M_t
 =\frac{y_Ay_B}{t-m_Y^2c^2}
 [\bar u_A(p_3)u_A(p_1)]
 [\bar u_B(p_4)u_B(p_2)].}
 \label{eq:yukawa-distinguishable-M-si-r10}
\end{equation}
两个双线性型各具有一阶动量量纲，与传播子分母的二阶动量量纲相消，所以 $\mathcal M_t$ 无量纲。

\begin{figure}[htbp]
 \centering
 \begin{tikzpicture}[x=1cm,y=1cm,>=Latex,baseline,line label/.style={fill=white,fill opacity=.96,text opacity=1,inner sep=1.2pt}]
  \coordinate (L1) at (-2.8,1.25);
  \coordinate (L2) at (-2.8,-1.25);
  \coordinate (V1) at (-0.75,0.45);
  \coordinate (V2) at (0.75,-0.45);
  \coordinate (R1) at (2.8,1.25);
  \coordinate (R2) at (2.8,-1.25);
  \draw[thick,->] (L1) -- (V1) node[line label,midway,above left=5pt] {$p_1$};
  \draw[thick,->] (V1) -- (R1) node[line label,midway,above right=5pt] {$p_1'$};
  \draw[thick,->] (L2) -- (V2) node[line label,midway,below left=5pt] {$p_2$};
  \draw[thick,->] (V2) -- (R2) node[line label,midway,below right=5pt] {$p_2'$};
  \draw[thick,dashed] (V1) -- (V2) node[line label,midway,above right=5pt] {$q$};
  \fill (V1) circle (1.4pt);
  \fill (V2) circle (1.4pt);
  \node at (0,-1.9) {$\displaystyle \frac{\ii}{q^2-m_Y^2+\ii0^+}$};
\end{tikzpicture}

 \caption{可区分 Dirac 费米子的树级标量交换。两个 Yukawa 顶角分别保持 $A$、$B$ 的费米子种类，中间标量传播子携带 $q=p_1-p_3$。}
 \label{fig:yukawa-exchange-r13}
\end{figure}

取等质量、等耦合 $m_A=m_B=m$、$y_A=y_B=y$。对两个初态自旋平均并对末态求和，得到
\begin{equation}
 \boxed{
 \overline{|\mathcal M_t|^2}
 =y^4\frac{(4m^2c^2-t)^2}{(t-m_Y^2c^2)^2}.}
 \label{eq:yukawa-distinguishable-M2-si-r10}
\end{equation}
质心系中若三动量大小为 $p$，则
\begin{equation*}
 s=4(m^2c^2+p^2),
 \qquad
 t=-2p^2(1-\cos\theta),
\end{equation*}
于是散射角分布为
\begin{equation}
 \boxed{
 \frac{\dd\sigma_Y}{\dd\Omega}
 =\frac{\hbar^2y^4}{64\pi^2s}
 \frac{[4m^2c^2+2p^2(1-\cos\theta)]^2}
 {[m_Y^2c^2+2p^2(1-\cos\theta)]^2}.}
 \label{eq:yukawa-distinguishable-angular-si-r10}
\end{equation}
当 $m_Yc\gg p$ 时，传播子在可访问动量转移范围内变化缓慢，角分布接近短程接触极限；当 $m_Yc\ll p$ 时，$t\simeq0$ 的前向区域增强。控制这两个极限的无量纲比值分别是 $p/(m_Yc)\ll1$ 与 $m_Yc/p\ll1$。

\begin{figure}[htbp]
 \centering
 \includegraphics[width=.78\linewidth]{yukawa_angular_distribution.png}
 \caption{可区分 Dirac 费米子经标量 Yukawa 媒介粒子弹性散射的归一化角分布。取 $mc/p=0.20$，比较 $m_Yc/p=0.30,1.0,3.0$；媒介粒子越轻，$t\simeq0$ 的前向增强越明显。}
 \label{fig:yukawa-angular-distribution}
\end{figure}

若四条外线属于同一种费米场，末态交换对应同一物理末态。费米子反对易关系给出 $t$、$u$ 两种收缩的相对负号，故
\begin{equation}
 \boxed{
 \mathcal M_{\rm id}
 =y^2\left[
 \frac{(\bar u_3u_1)(\bar u_4u_2)}{t-m_Y^2c^2}
 -
 \frac{(\bar u_4u_1)(\bar u_3u_2)}{u-m_Y^2c^2}
 \right].}
 \label{eq:yukawa-identical-M-si-r10}
\end{equation}
两种通道的干涉由无量纲平方振幅中的标量函数
\begin{equation}
 \boxed{
 \mathcal I_{tu}
 =\frac12\left(4m^2c^2s-4m^2c^2t+st+t^2\right)}
 \label{eq:yukawa-interference-invariants-si-r10}
\end{equation}
控制。未偏振平方因此为
\begin{align}
 \overline{|\mathcal M_{\rm id}|^2}
 =&\frac{y^4(4m^2c^2-t)^2}{(t-m_Y^2c^2)^2}
 +\frac{y^4(4m^2c^2-u)^2}{(u-m_Y^2c^2)^2}\notag\\
 &-\frac{2y^4\mathcal I_{tu}}
 {(t-m_Y^2c^2)(u-m_Y^2c^2)}.
 \label{eq:yukawa-identical-M2-general-si-r10}
\end{align}
若在全 $4\pi$ 上对有序末态变量积分，还需乘全同粒子的 $1/2!$ 对称因子。

低速极限把相对论交换振幅连接到熟悉的 Yukawa 势。控制条件为
\begin{equation*}
 \frac{|\mathbf p|}{mc}\ll1,
 \qquad
 \frac{|\mathbf p'|}{mc}\ll1,
 \qquad
 \frac{|q^0|}{|\mathbf q|}\ll1,
\end{equation*}
此时 $t\simeq-\mathbf q^2$，可区分费米子的主导振幅为
\begin{equation}
 \mathcal M_t
 \simeq-\frac{4m^2c^2y^2}{\mathbf q^2+m_Y^2c^2}.
 \label{eq:yukawa-NR-M-si-r10}
\end{equation}
把\eqnrefs{eq:yukawa-NR-M-si-r10} 代入相对论--非相对论匹配\eqnrefs{eq:relativistic-to-nr-potential-matching-dp75}，沿用其中已经定义的统一归一化。对两条外线质量均为 $m$ 的当前通道，立即得到
\begin{equation}
 \widetilde V_Y(\mathbf q)
 =-\frac{4\pi\alpha_Y\hbar^3c}
 {\mathbf q^2+m_Y^2c^2},
 \qquad
 \alpha_Y\equiv\frac{y^2}{4\pi},
 \label{eq:yukawa-Vq-si-r10}
\end{equation}
其三维 Fourier 变换给出
\begin{equation}
 \boxed{
 V_Y(r)
 =-\alpha_Y\hbar c\,\frac{\ee^{-r/\lambda_Y}}{r},
 \qquad
 \lambda_Y=\frac{\hbar}{m_Yc}.}
 \label{eq:yukawa-potential-SI-r10}
\end{equation}

\begin{exbox}[媒介质量怎样直接变成作用程]
把上面的精确结果用于一个具体尺度。若媒介粒子的静能为
\begin{equation*}
 m_Yc^2=80\,\mathrm{GeV},
\end{equation*}
则
\begin{equation*}
 \lambda_Y
 =\frac{\hbar c}{m_Yc^2}
 \simeq\frac{0.19732698\,\mathrm{GeV\,fm}}{80\,\mathrm{GeV}}
 \simeq2.47\times10^{-3}\,\mathrm{fm}
 =2.47\times10^{-18}\,\mathrm m.
\end{equation*}
因此在 $r=\lambda_Y$ 处，势相对于同耦合的无质量 $1/r$ 核只剩 $\ee^{-1}\simeq0.368$；在 $r=5\lambda_Y$ 时只剩
\begin{equation*}
 \ee^{-5}\simeq6.74\times10^{-3}.
\end{equation*}
反过来，若媒介无质量，$\lambda_Y\to\infty$，任何有限 $r$ 都不会得到指数压低。这一数值算例把“重媒介对应短程、无质量媒介对应长程”从一句定性描述变成了可以直接比较的长度尺度。
\end{exbox}

这个势是树级低能匹配得到的有效相互作用。若进一步只用 Schr\"odinger 第一 Born 项计算散射，还需在实际入射能量与部分波工作空间 $\mathcal W$ 上满足
\begin{equation*}
 \epsilon_{\rm Born,Y}
 \equiv
 \|G_0^{(+)}(E)V_Y\|_{\mathcal W}\ll1,
\end{equation*}
此时多重势插入从 $O(\epsilon_{\rm Born,Y}^2)$ 开始修正一次插入结果。传播子极点、有限作用程和前向增强都由同一个分母 $q^2-m_Y^2c^2$ 控制；当 $\epsilon_{\rm Born,Y}$ 不再小，应保留该势但求解完整 Lippmann--Schwinger 方程，而不是继续截断 Born 级数。



对过程 $\psi(p_1)+\bar\psi(p_2)\to\phi(k_1)+\phi(k_2)$，两个全同末态标量允许两种连接次序。树级约化振幅为
\begin{align}
 \mathcal M
 =y^2\bar v(p_2)\Bigg[
 &\frac{\slashed p_1-\slashed k_1+mc}
 {(p_1-k_1)^2-m^2c^2+\ii0^+}
 +\frac{\slashed p_1-\slashed k_2+mc}
 {(p_1-k_2)^2-m^2c^2+\ii0^+}
 \Bigg]u(p_1).
 \label{eq:yukawa-annihilation-two-scalar-si-r10}
\end{align}
阈值条件是 $s\ge4m_Y^2c^2$，且二者必须在振幅层先相加再平方。该拓扑与 Compton 散射共享“同一内部费米线的两种玻色子连接次序”这一结构。


\begin{exbox}[赝标量耦合的低能自旋结构]
若媒介粒子是赝标量，采用厄米相互作用
\begin{equation}
 \Lag_P=-y_P\sqrt{\hbar c}\,\phi_P\bar\psi\ii\gamma^5\psi.
 \label{eq:pseudoscalar-yukawa-si-r10}
\end{equation}
在 $|\mathbf p|/(mc),|\mathbf p'|/(mc)\ll1$ 下，赝标量双线性型的首项是
\begin{equation}
 \boxed{
 \bar u(p')\ii\gamma^5u(p)
 =-\ii\,\xi'^\dagger
 \boldsymbol\sigma\cdot(\mathbf p'-\mathbf p)\xi
 +\order\!\left(\frac{p^3}{m^2c^2}\right).}
 \label{eq:pseudoscalar-NR-bilinear-si-r10}
\end{equation}
因此两个赝标量顶角产生 $[\boldsymbol\sigma_1\cdot\mathbf q][\boldsymbol\sigma_2\cdot\mathbf q]$ 的分子，坐标空间势分解为自旋--自旋项与张量项。媒介粒子质量决定作用程，而顶角的 Lorentz 结构决定低能自旋依赖。


\medskip

赝标量耦合与标量耦合的低能行为不同，原因直接来自在壳旋量之间的 $\bar u\ii\gamma^5u$。从正文 Yukawa 相互作用~\eqnrefs{eq:yukawa-full-si-r10} 选出赝标量顶角，把正能 Dirac 旋量按 $|\mathbf p|/(mc)\ll1$ 展开到首个非零阶，可以看到常数项相消、结果与动量转移成正比；这正是\eqnrefs{eq:pseudoscalar-NR-bilinear-si-r10}。

在 Dirac 表象中
\begin{equation*}
 \gamma^5=
 \begin{pmatrix}0&I\\ I&0\end{pmatrix},
 \qquad
 \bar u=u^\dagger\gamma^0.
\end{equation*}
代入协变归一化旋量，先保持精确能量：
\begin{align*}
 \bar u(p')\ii\gamma^5u(p)
 &=\ii\sqrt{\frac{(E'+mc^2)(E+mc^2)}{c^2}}\,
 \xi'^\dagger\\
 &\quad\times\left[
 \frac{c\boldsymbol\sigma\cdot\mathbf p}{E+mc^2}
 -\frac{c\boldsymbol\sigma\cdot\mathbf p'}{E'+mc^2}
 \right]\xi.
\end{align*}
低速展开给
\begin{equation*}
 \sqrt{\frac{(E'+mc^2)(E+mc^2)}{c^2}}
 =2mc+\order\!\left(\frac{p^2}{mc}\right),
\end{equation*}
而
\begin{equation*}
 \frac{c}{E+mc^2}
 =\frac{1}{2mc}\left[1+\order\!\left(\frac{p^2}{m^2c^2}\right)\right].
\end{equation*}
两者相乘后主项为
\begin{equation*}
 \bar u(p')\ii\gamma^5u(p)
 =-\ii\,\xi'^\dagger\boldsymbol\sigma\cdot(\mathbf p'-\mathbf p)\xi
 +\order\!\left(\frac{p^3}{m^2c^2}\right),
\end{equation*}
即目标式。与标量双线性型的常数主项相比，赝标量顶角在低能处多出一个动量转移，因此产生自旋依赖而不是纯中心势。
\end{exbox}
